\documentclass[11pt]{article}

\usepackage[margin=1in]{geometry}
\usepackage{amsmath,amssymb,amsthm}
\usepackage{mathtools}
\usepackage{booktabs}
\usepackage{array}
\usepackage{enumitem}
\usepackage{hyperref}
\usepackage{xcolor}

\hypersetup{
  colorlinks=true,
  linkcolor=blue!50!black,
  urlcolor=blue!50!black,
  citecolor=blue!50!black
}

\title{Problem Formulation:\\
Empirical Characterization of the Feasible Set $F$\\
for Contact-Rich Manipulation Primitives}
\author{NAMO Research Notes}
\date{}

\begin{document}
\maketitle

\begin{abstract}
Paper-style definitions of every object, operation, and claim in the $F$
characterization work. Intended as the spine for the ``Preliminaries'' /
``Problem Formulation'' section of the paper, and as the reference any later
document or experiment can point back to.
\end{abstract}

\section{Setting}

We study \textbf{contact-rich manipulation with parameterized primitives}. The
robot operates in a planar environment containing static obstacles (walls),
movable obstacles (boxes), and a navigation goal. The robot executes
parameterized push primitives on movable obstacles to change scene
configuration and enable navigation. We instantiate this setting in
\textbf{NAMO} (Navigation Among Movable Obstacles) and characterize the
structure of primitive feasibility empirically.

\section{Positioning: Controller-Aware Planning and the Action-Space Narrow Passage}

This work occupies a specific position in the manipulation research
landscape. We make it explicit before introducing formal machinery, because
the framing determines what the contribution is --- and what it is not.

\subsection{Two camps: robust controller design vs.\ controller-aware planning}

Manipulation research divides along a methodological axis with two poles.

\paragraph{Pole 1 --- Robust controller design.}
The intellectual effort is concentrated in the controller. The goal is to
engineer (or learn) a controller whose success region across scenes is so
large that planning around it becomes unnecessary. Examples: classical
compliant control (Mason, Whitney, Salisbury), large-scale visuomotor
policies (RT-X, $\pi$, OpenVLA), domain-randomized end-to-end training. The
bet is that controllers can be made robust enough that explicit reasoning
about their failure regions is dispensable.

\paragraph{Pole 2 --- Controller-aware planning.}
The intellectual effort is concentrated in the planner. The controller is
taken as given; the planner reasons explicitly about where it succeeds and
fails. Examples: funnel composition (Burridge--Rizzi--Koditschek 1999),
LQR-Trees (Tedrake et al.\ 2010), contact trust regions (Suh et al.\ 2023),
TAMP with feasibility samplers (Garrett, Kaelbling, Lozano-P\'erez). The bet
is that controllers always have non-trivial failure regions, that those
regions have structure worth understanding, and that planning over a
characterized success set produces system-level capability beyond what
improving the controller alone can achieve.

\paragraph{This work is in Pole 2.}
We take a contact-rich push primitive as given, characterize its feasible set
$F$ empirically, and study how the structure of $F$ should inform models
that sample within it. We do not contest Pole 1's bet; we observe that
controller-aware planning is an underexplored axis in contact-rich
manipulation that retains value regardless of how robust controllers become.

\subsection{The narrow-passage problem in action space}

Sampling-based motion planning has a half-century-old lineage devoted to the
\textbf{narrow-passage problem}: configuration spaces where success requires
sampling within a thin region of $\mathcal{C}_{\text{free}}$, and uniform
sampling fails because the passage's measure is small (Hsu, Latombe, Motwani
1997). Bridge sampling (Hsu et al.), Gaussian-near-obstacles sampling (Boor
et al.), and learned samplers (Ichter et al.\ 2018) are the classical
responses: characterize the passage geometry, then design samplers that put
mass on it.

Our work is the \textbf{action-space analogue} of the narrow-passage problem.
The correspondence is structural:

\begin{center}
\renewcommand{\arraystretch}{1.15}
\begin{tabular}{ll}
\toprule
Classical motion planning & This work \\
\midrule
Configuration space $\mathcal{C}$ & Action-parameter space $\mathcal{A}$ \\
Free space $\mathcal{C}_{\text{free}}$ & Reachable set $R(s)$ \\
Goal-reaching subspace & Feasible set $F(s, g)$ \\
Narrow passage in $\mathcal{C}_{\text{free}}$ & Sparse $F$ within large $R$ --- low $\rho = |F|/|R|$ \\
Sample complexity $\sim 1/\epsilon^d$ & Expected uniform-sampler trials $\sim 1/\rho$ \\
Bridge / Gaussian / learned samplers & Structure-informed classifiers and generative samplers \\
\bottomrule
\end{tabular}
\end{center}

The difficulty metric $\rho = |F|/|R|$ adopted in
Section~\ref{sec:difficulty} is the discrete-action-space analogue of the
$\epsilon$-good measure from the narrow-passage literature. The
structure-informed model variants (Tier~3 baselines,
Section~\ref{sec:baselines}) are the action-space analogues of bridge
sampling --- non-uniform sampling distributions designed to place mass on the
success region.

Two differences with the classical setting are essential, and they delimit
what is genuinely new:
\begin{enumerate}[leftmargin=2em,itemsep=2pt]
  \item \textbf{The passage is in action space, not configuration space.} The
        narrow region is somewhere the planner has to \emph{aim}, not
        somewhere the robot has to \emph{thread through}. Local geometry near
        $F$ is a metric on $\mathcal{A}$, governed by the dynamics $T$ and
        the success-state set $G(g)$, not by configuration-space obstacles.
  \item \textbf{The passage is dynamics-induced, not geometry-induced.} $F$'s
        narrowness is determined by contact-rich dynamics: which pushes
        happen to land the object in the small set of poses that satisfy the
        goal. Classical narrow passages are facts about obstacle layout. Our
        narrow passages are facts about how $T$ maps actions to states and
        how that state distribution intersects $G$. The structural
        characterization presented here (face $\to$ contact $\to$ depth
        bottleneck, wall dependence, depth-window contiguity) is empirical
        knowledge about the dynamics-induced passage geometry that has no
        analytical analogue in the classical literature.
\end{enumerate}

The recursive multi-push extension ($F_k'$, Section~\ref{sec:npush}) has no
clean classical analogue. The closest classical setting --- sequential
narrow passages in motion planning --- has not been characterized empirically
for either configuration-space or action-space cases. Our recursive
characterization is therefore novel by analogy as well as in substance.

\subsection{What this positioning implies}

These two framings --- controller-aware planning and the action-space narrow
passage --- together specify the contribution:
\begin{itemize}[leftmargin=2em,itemsep=0pt]
  \item The work belongs to the controller-aware planning lineage (funnel
        composition, LQR-Trees, contact trust regions), extended to
        contact-rich manipulation in clutter where analytical RoA computation
        is intractable.
  \item It transfers the methodological discipline of the narrow-passage
        literature (characterize the passage, then design the sampler) from
        configuration-space motion planning to action-space manipulation
        primitives.
  \item It contrasts cleanly with the dominant Pole~1 trajectory in current
        manipulation research --- not as a competitor but as a complementary
        axis whose value is independent of how robust controllers become.
\end{itemize}

The baseline design (Section~\ref{sec:baselines}) operationalizes this
positioning: geometric-heuristic baselines (Tier~1) are the action-space
analogues of bridge sampling; learned baselines (Tier~2) are the analogues of
Ichter-style learned samplers; structure-informed baselines (Tier~3) test
whether the empirical characterization causally improves the sampler. Showing
the layered lift across tiers is the action-space replication of the
narrative the narrow-passage literature established for configuration-space
motion planning.

\section{State and Conditioning}

Let $s \in \mathcal{S}$ denote a \textbf{scene state}, comprising:
\begin{itemize}[leftmargin=2em,itemsep=0pt]
  \item Static obstacle geometry (walls).
  \item Movable obstacle poses $\{(x_i, y_i, \theta_i)\}_{i=1}^M$.
  \item Robot pose $q_r = (x_r, y_r, \theta_r)$.
\end{itemize}

Let $g \in \mathcal{G}$ denote a \textbf{task goal}, here a target robot pose
$g = (x_g, y_g, \theta_g)$. The conditioning context for the feasibility
analysis is the tuple $(s, g)$.

\section{Action Space}

Let $\mathcal{A}$ denote the \textbf{primitive action space}: the parameter
space of the push skill. In our NAMO instantiation,
\[
  \mathcal{A} \;=\; \mathcal{O} \times \mathcal{F} \times \mathcal{P} \times \mathcal{D}
\]
where $\mathcal{O}$ is the chosen movable obstacle,
$\mathcal{F} = \{1,2,3,4\}$ is the face index,
$\mathcal{P} = \{1, \dots, 15\}$ is the contact-point index along the face,
and $\mathcal{D} = \{1, \dots, 10\}$ is the discretized push depth. For a
fixed object, $|\mathcal{A}| = 600$.

We treat $\mathcal{A}$ as a \textbf{discretization of a continuous parameter
space} $\mathcal{A}_c$ in (push direction, push depth), and verify that key
structural properties are stable across discretization resolutions.

\section{Reachability}

The \textbf{reachability set} $R(s) \subseteq \mathcal{A}$ is the subset of
primitives the robot can physically execute from state $s$:
\[
  R(s) \;=\; \{\,a \in \mathcal{A} \;:\;
    \text{the robot can reach the contact point of $a$ in $s$}\,\}.
\]
In NAMO, $R(s)$ is computed exactly and at low cost via wavefront BFS over an
inflated occupancy grid. Reachability is \textbf{task-independent}: it depends
on $s$ but not on $g$. In the narrow-passage analogy of Section~2, $R(s)$
plays the role of free space $\mathcal{C}_{\text{free}}$ --- the ambient
region within which the success region $F$ is embedded.

\section{Push Dynamics}

Let $T : \mathcal{S} \times \mathcal{A} \rightharpoonup \mathcal{S}$ denote
the \textbf{push transition operator}: $T(s, a)$ is the post-push scene state
obtained by executing primitive $a$ from state $s$ in simulation. $T$ is
partial (defined only for $a \in R(s)$), discontinuous (contact-mode
switches), and not analytically tractable. The intractability of $T$ is what
forces empirical characterization of $F$ --- analytical preimage computation
is unavailable.

\section{Success Criterion}

Given a task goal $g$, the \textbf{success-state set} is
\[
  G(g) \;=\; \{\,s' \in \mathcal{S} \;:\; g \text{ is reachable for the robot from } s'\,\}.
\]
Equivalently, $G(g)$ is the set of post-push states from which a navigation
plan to $g$ exists. Membership in $G(g)$ is decided by a wavefront
reachability query --- exact, fast, no learning involved.

\section{The Feasible Set $F$}

The central object of study is the \textbf{feasible set}:
\[
  F(s, g) \;=\; \{\,a \in R(s) \;:\; T(s, a) \in G(g)\,\}.
\]

Equivalently, $F$ is the \textbf{preimage} in action space of the
success-state set $G(g)$ under the push dynamics $T(s, \cdot)$, restricted to
reachable actions:
\[
  F(s, g) \;=\; R(s) \,\cap\, T(s, \cdot)^{-1}\!\big(G(g)\big).
\]

We adopt the decomposition $F = R \cap C$, where
$C(s, g) = \{\,a \in \mathcal{A} \,:\, T(s, a) \in G(g)\,\}$ is the
\textbf{clearance set} (task success conditional on execution). $R$ is exact
and free; $C$ is what a learned model must capture.

\paragraph{Conceptual framing.}
$F(s, g)$ is the \emph{empirical region of attraction in action-parameter
space}, conditioned on scene and task goal. It is the manipulation analogue
of classical control-theoretic regions of attraction (Lyapunov, SOS, contact
trust regions), but defined in action space rather than state space and
characterized empirically rather than analytically --- because contact-rich
dynamics admit no closed-form preimage. Through the lens of Section~2, $F$
is the \textbf{narrow passage in action space} induced by the controller's
contact-rich dynamics.

\section{Difficulty}\label{sec:difficulty}

We define instance \textbf{difficulty} operationally as
\[
  \rho(s, g) \;=\; \frac{|F(s, g)|}{|R(s)|}.
\]
$\rho$ measures the sparsity of feasible primitives within the reachable set.
Its inverse $1/\rho$ is the expected number of uniform random primitive
evaluations required to find a feasible action --- the discrete-action-space
analogue of the $\epsilon$-narrowness measure from Hsu--Latombe--Motwani
(1997). We stratify instances into difficulty buckets by $\rho$ (Very Hard,
Hard, Medium, Easy, Very Easy) and report all structural results per stratum.

\section{Structural Descriptors of $F$}\label{sec:descriptors}

For each instance $(s, g)$ with $|F(s, g)| > 0$, we compute structural
descriptors over the discrete action grid:
\begin{itemize}[leftmargin=2em,itemsep=0pt]
  \item \textbf{Active face count} $n_F \in \{0,\dots,4\}$: number of faces
        containing any feasible primitive.
  \item \textbf{Contact-point span} $n_P$: number of active contact points
        within active faces.
  \item \textbf{Depth window} $[d_{\min}, d_{\max}]$ per active contact point
        and its \textbf{contiguity} (gap fraction).
  \item \textbf{Connected-component count} $K(s, g)$: number of disconnected
        feasible regions in the (face, contact, depth) grid under face-aware
        adjacency.
  \item \textbf{Wall-collision involvement} $w(s, g)$: fraction of
        $a \in F(s, g)$ for which simulator rollout reports object-wall
        contact during execution.
\end{itemize}

These descriptors operationalize the \emph{structure} of $F$. In the
narrow-passage analogy, they are the action-space analogues of measurements
one would make on a configuration-space passage: how thin, how connected, how
aligned with which obstacle features.

\section{Empirical Characterization Protocol}

For each scene $s$ in a generated dataset $\mathcal{D}_{\text{scenes}}$ and
corresponding task goal $g$:
\begin{enumerate}[leftmargin=2em,itemsep=0pt]
  \item Compute $R(s)$ via wavefront BFS.
  \item For each $a \in R(s)$, execute $T(s, a)$ in simulation.
  \item Evaluate $T(s, a) \in G(g)$ via post-push wavefront reachability.
  \item Record per-primitive outcome (success / failure), collision flags,
        and post-push pose.
  \item Compute structural descriptors of $F(s, g)$.
\end{enumerate}

We refer to this as \textbf{exhaustive 1-push $F$ characterization}.

\section{The Learned Predictor}

We are interested in models $p_\theta(a \mid s, g)$ that \textbf{predict} or
\textbf{sample} primitives intended to lie in $F(s, g)$. We consider two
families:
\begin{itemize}[leftmargin=2em,itemsep=0pt]
  \item \textbf{Classifier}: $h_\theta(s, g, a) \in [0, 1]$ approximating
        $\Pr[\,a \in F(s, g)\,]$, trained on per-primitive labels. Sampling
        proceeds by ranking $a \in R(s)$ by $h_\theta$ and drawing top-$k$ or
        threshold-filtered samples.
  \item \textbf{Generative sampler} (diffusion): $p_\theta(a \mid s, g)$
        trained on $(s, g, a)$ tuples with $a \in F(s, g)$. Sampling proceeds
        by reverse diffusion conditioned on $(s, g)$.
\end{itemize}

For both families, we say the model \textbf{implicitly characterizes $F$} if
its predictions respect the structural descriptors of $F$. Both families are
instances of the controller-aware planning paradigm of Section~2: the
underlying push controller is held fixed; the model learns where to invoke
it.

\section{Baselines}\label{sec:baselines}

We organize baselines into tiers, each isolating a distinct claim. Every
baseline tests one specific component of the hypothesis chain.

\paragraph{Tier 0 --- Bounds.}
\begin{itemize}[leftmargin=2em,itemsep=0pt]
  \item \textbf{B0a --- Oracle}: exhaustive evaluation of $R(s)$, returning
        ground-truth $F(s,g)$. Upper bound.
  \item \textbf{B0b --- Random over $\mathcal{A}$}: uniform sampling
        ignoring reachability. Trivial floor.
  \item \textbf{B0c --- Random over $R(s)$}: uniform sampling restricted to
        reachable primitives. Achieves precision $= \rho$ in expectation.
        \emph{The critical floor:} the gap between any method and B0c
        isolates what the method has learned about clearance $C$ beyond the
        free reachability oracle.
\end{itemize}

\paragraph{Tier 1 --- Geometric heuristics (action-space analogues of bridge sampling).}
\begin{itemize}[leftmargin=2em,itemsep=0pt]
  \item \textbf{B1a --- Push-toward-open-space}: bias samples by the
        wavefront-distance gradient at the contact point.
  \item \textbf{B1b --- Wall-aware sampling}: bias samples toward primitives
        that produce object-wall contact.
  \item \textbf{B1c --- Center-of-face heuristic}: push from the center of
        an object face at moderate depth.
  \item \textbf{B1d --- Hierarchy-respecting random}: sample face $\to$
        contact-point $\to$ depth uniformly within active levels.
\end{itemize}

\paragraph{Tier 2 --- Simple learned baselines.}
\begin{itemize}[leftmargin=2em,itemsep=0pt]
  \item \textbf{B2a --- Logistic regression} on hand-engineered features.
  \item \textbf{B2b --- Random forest} on the same features.
  \item \textbf{B2c --- Independent-logits classifier}: deep network
        outputting 600 independent logits, structure-blind.
  \item \textbf{B2d --- Independent-logits + rejection over $R$}: B2c with
        reachability filtering at inference. The cleanest classifier baseline
        for the generative comparison.
\end{itemize}

\paragraph{Tier 3 --- Structure-informed learned baselines (the contributions).}
Each is an architectural variant of B2d that exploits one structural finding
from Section~\ref{sec:descriptors}.
\begin{itemize}[leftmargin=2em,itemsep=0pt]
  \item \textbf{B3a --- Conv-over-grid head}: convolutions over the
        (contact-point, depth) grid per face. Tests the contiguity prior.
  \item \textbf{B3b --- Hierarchical face $\to$ contact $\to$ depth head}:
        sequential prediction respecting the bottleneck hierarchy.
  \item \textbf{B3c --- Wall-distance-augmented input}: scene encoder
        receives a wall-distance field around the object.
\end{itemize}

\paragraph{Tier 4 --- The generative comparison.}
\begin{itemize}[leftmargin=2em,itemsep=0pt]
  \item \textbf{B4a --- Diffusion model}: scene + goal conditioned, samples
        primitives, restricted to $R$ at inference.
  \item \textbf{B4b --- Diffusion ablation, scene-only}: diagnostic for the
        importance of goal conditioning.
\end{itemize}

\paragraph{Reporting protocol.}
Every baseline reports precision@$k$, coverage@$k$, and structural-alignment
metrics per difficulty stratum. For each baseline we state in advance --- based
on the characterization of $F$ --- how it should perform; the agreement or
disagreement is itself a hypothesis test on the structural findings.

\paragraph{Minimum viable baseline set.}
If time-constrained, the minimum baselines that license the paper's central
claims are: \textbf{B0c} (uniform over $R$), \textbf{B1b} (wall-aware),
\textbf{B2d} (deep classifier + rejection), \textbf{B3b} (hierarchical head),
\textbf{B4a} (diffusion). Five baselines, each licensed by a specific
component of the hypothesis chain.

\section{Evaluation Against Ground-Truth $F$}

Given ground-truth $F(s, g)$ for each test instance and $k$ samples
$\{a_1, \dots, a_k\} \sim p_\theta(\cdot \mid s, g)$, we report:
\begin{itemize}[leftmargin=2em,itemsep=0pt]
  \item \textbf{Precision@$k$}:
        $\frac{1}{k} \sum_{i=1}^{k} \mathbf{1}[\,a_i \in F(s, g)\,]$.
  \item \textbf{Coverage@$k$}: fraction of connected components of $F(s, g)$
        hit by at least one sample.
  \item \textbf{Distributional divergence}: JS divergence between the
        empirical sample distribution and the uniform distribution over
        $F(s, g)$.
  \item \textbf{Structural alignment}: per-difficulty-stratum precision and
        coverage; correlation between failure mode and instance descriptors.
\end{itemize}

The \textbf{lift over a uniform-over-$R$ baseline} (B0c) isolates whether
the model has learned anything about $C$ beyond reachability.

\section{Hypotheses}

We pre-state hypotheses tested by the characterization (these were registered
before data collection):
\begin{itemize}[leftmargin=2em,itemsep=2pt]
  \item \textbf{H1 (Difficulty)}: Instance difficulty is governed by
        $\rho = |F|/|R|$, not $|F|$ alone.
  \item \textbf{H2 (Contiguity)}: $F$ is spatially contiguous in
        $\mathcal{A}$, not scattered.
  \item \textbf{H3 (Fragmentation)}: Hard instances have higher fragmentation
        $K(s, g)$ than easy instances. \emph{(Falsified by data: hard
        instances are predominantly unimodal.)}
  \item \textbf{H4 (Depth precision)}: On hard instances, $F$ requires
        precise depth but spans multiple contact points. \emph{(Refined by
        data: contact-point precision is the dominant bottleneck; depth
        windows are tight but contiguous.)}
  \item \textbf{H5 (Wall dependence)}: On hard instances, a disproportionate
        fraction of $F$ involves wall collisions. \emph{(Confirmed: 76\% on
        Very Hard.)}
\end{itemize}

\section{The $N$-Push Extension (Defined for Completeness; Deferred)}\label{sec:npush}

For $n$-push problems where $F(s, g) = \emptyset$, we define recursively:
\[
  F_k'(s, g) \;=\; \{\,a \in R(s) \,:\, F_{k-1}'\big(T(s, a),\, g\big) \neq \emptyset\,\},
  \qquad F_1'(s, g) \;\coloneqq\; F(s, g).
\]
$F_k'$ is the set of first-of-$k$ pushes that lead to a state from which a
$(k-1)$-push solution exists. This is \textbf{recursive feasibility
composition} --- feasibility flows backward from the goal through the chain
--- and is the action-space analogue of funnel composition
(Burridge--Rizzi--Koditschek) and LQR-Trees (Tedrake).

The structural questions about $F$ extend recursively: does $F_k'$ inherit
contiguity, unimodality, and wall dependence from $F_1'$? The closest
classical analogue --- sequential narrow passages in motion planning --- has
not been characterized empirically for either configuration-space or
action-space cases. This forms the basis of the multi-push study and is left
to future work.

\section{Scope and Disclaimers}

\begin{itemize}[leftmargin=2em,itemsep=0pt]
  \item We do not claim $F$'s structure transfers to manipulation primitives
        outside push-in-clutter without further characterization.
  \item We do not claim sim-measured $F$ matches real-world $F$; sim-to-real
        validation is left to future work.
  \item We do not claim our learned models are state-of-the-art on any
        external benchmark; the contribution is structural, not numerical.
  \item We do not claim the methodology is novel in concept; empirical
        characterization of feasibility regions has precedent (Dex-Net for
        grasping, contact trust regions for few-contact systems). We claim
        novelty in \emph{application} (contact-rich pushing in clutter), in
        \emph{positioning} (controller-aware planning + action-space
        narrow-passage framing), and in \emph{coupling} characterization to
        classifier-first methodology with falsifiable hypotheses.
  \item We do not contest the robust-controller-design (Pole~1) bet of
        Section~2. Controller-aware planning is complementary to controller
        improvement; both can be pursued in parallel.
\end{itemize}

\section*{Notation Summary}

\begin{center}
\renewcommand{\arraystretch}{1.2}
\begin{tabular}{>{$}l<{$} l}
\toprule
\textnormal{Symbol} & Meaning \\
\midrule
s                              & Scene state (geometry + robot pose) \\
g                              & Task goal (target robot pose) \\
\mathcal{A}                    & Primitive action space (600 push primitives in NAMO) \\
R(s)                           & Reachable subset of $\mathcal{A}$ --- exact via wavefront BFS \\
T(s, a)                        & Post-push scene state \\
G(g)                           & Success-state set (states from which $g$ is reachable) \\
C(s, g)                        & Clearance set $\{a : T(s, a) \in G(g)\}$ \\
F(s, g)                        & Feasible set $R(s) \cap C(s, g)$ \\
\rho(s, g)                     & Difficulty $|F| / |R|$ \\
K(s, g)                        & Number of connected components of $F$ \\
w(s, g)                        & Wall-collision fraction in $F$ \\
p_\theta(a \mid s, g)          & Learned sampler \\
h_\theta(s, g, a)              & Learned classifier \\
\bottomrule
\end{tabular}
\end{center}

\section*{Why This Document Exists}

A paper lives or dies on whether the reader can hold the central object in
their head precisely by the end of Section~2. With these definitions in
place, every later claim (``$F$ is unimodal on hard instances,'' ``the model
fails on wall-dependent $F$,'' ``diffusion does not lift over
classifier+rejection on hard cases'') becomes a precise statement about a
precisely-defined object, evaluated by a precisely-defined protocol. That
precision is what separates a paper that gets engaged with from a paper that
gets called ``interesting but unclear.''

Pin this. Refer back to it. When confused about what an experiment tests,
reread the definitions. They are the spine.

\end{document}
